% GENERATED FILE. Edit canonical lessons, then run npm run generate:books.
% canonical-source-hash: fnv1a64:2c2acae47a4450a5
% canonical-lessons: RU-C474-nervnichat, RU-C474-trevozhnyy, RU-C474-bespokoitsya, RU-C474-bespokoystvo, RU-C474-uspokaivatsya

\chapter{Worry and Calm}
\label{ch:ru-worry-and-calm}
\hlchaptermodality{474}

\begin{chapteropening}
\textbf{By the end of this chapter:} \emph{I can talk about nerves, worry and calming down.}

Complete the last lesson of chapter 474: I can talk about nerves, worry and calming down.
\end{chapteropening}

\section[\texorpdfstring{\ru{нервничать} (nérvnichat)}{нервничать (nérvnichat)}]{\ru{нервничать} (nérvnichat) — to be nervous}

\label{lesson:RU-C474-nervnichat}



\begin{quote}
Before the new one: say the Russian for despair, then the Russian for to sigh.
\end{quote}

\subsection*{You'll want to know: \ru{нервничать}}
\textbf{\ru{нервничать}} (\emph{nérvnichat}) — \textquotedblleft{}to be nervous\textquotedblright{}.

\begin{quote}
\ru{Перед} \ru{экзаменом} \ru{я} \ru{всегда} \ru{нервничаю}.
\end{quote}

\begin{quote}
\emph{Péred ekzámenom ya vsegdá nérvnichayu.}
\end{quote}

\begin{quote}
I always get nervous before an exam.
\end{quote}

Imperfective; present \ru{нервничаю}. Perfective \ru{понервничать} (for a while).

\begin{cousinweb}
\emph{\ru{нервничать}} is built on \emph{\ru{нервный}} 'nervous', from \emph{\ru{нерв}} 'nerve', which EtymDB traces to Latin \emph{nervus} 'sinew' (Lewis \& Short: a sinew, a tendon).
\end{cousinweb}

\begin{grammarlens}[title={What you've built}]
One more word for telling, explaining and describing at B1.
\end{grammarlens}

\subsection*{Your turn}
\begin{itemize}
\raggedright
  \item \emph{Say it:} \emph{nérvnichat}
  \item \emph{Say it:} \emph{nérvnichat}, once more
  \item \emph{Say it:} a sentence of your own with \emph{nérvnichat}
  \item \emph{Recall:} say the Russian for despair, then the Russian for to sigh, then say \emph{nérvnichat} again
\end{itemize}

\begin{tcolorbox}[breakable,colback=teal!4,colframe=teal!35!black,title={Before you move on}]
What does \ru{нервничать} mean? (\textquotedblleft{}To be nervous\textquotedblright{}.) Say it once more.
\end{tcolorbox}
\section[\texorpdfstring{\ru{тревожный} (trevózhnyy)}{тревожный (trevózhnyy)}]{\ru{тревожный} (trevózhnyy) — anxious; alarming}

\label{lesson:RU-C474-trevozhnyy}



\begin{quote}
Before the new one: say the Russian for to sigh, then the Russian for to be nervous.
\end{quote}

\subsection*{You'll want to know: \ru{тревожный}}
\textbf{\ru{тревожный}} (\emph{trevózhnyy}) — \textquotedblleft{}anxious; alarming\textquotedblright{}.

\begin{quote}
\ru{У} \ru{неё} \ru{был} \ru{тревожный} \ru{голос}.
\end{quote}

\begin{quote}
\emph{U neyó byl trevózhnyy gólos.}
\end{quote}

\begin{quote}
Her voice sounded anxious.
\end{quote}

Of a person, anxious; of news, alarming.

\begin{cousinweb}
\emph{\ru{тревожный}} is built on \emph{\ru{тревога}} 'alarm, anxiety', which is inherited from Proto-Slavic \emph{tr\ru{ь}voga} (EtymDB).
\end{cousinweb}

\begin{grammarlens}[title={What you've built}]
One more word for telling, explaining and describing at B1.
\end{grammarlens}

\subsection*{Your turn}
\begin{itemize}
\raggedright
  \item \emph{Say it:} \emph{trevózhnyy}
  \item \emph{Say it:} \emph{trevózhnyy}, once more
  \item \emph{Say it:} a sentence of your own with \emph{trevózhnyy}
  \item \emph{Recall:} say the Russian for to sigh, then the Russian for to be nervous, then say \emph{trevózhnyy} again
\end{itemize}

\begin{tcolorbox}[breakable,colback=teal!4,colframe=teal!35!black,title={Before you move on}]
What does \ru{тревожный} mean? (\textquotedblleft{}Anxious; alarming\textquotedblright{}.) Say it once more.
\end{tcolorbox}
\section[\texorpdfstring{\ru{беспокоиться} (bespokóitsya)}{беспокоиться (bespokóitsya)}]{\ru{беспокоиться} (bespokóitsya) — to worry}

\label{lesson:RU-C474-bespokoitsya}



\begin{quote}
Before the new one: say the Russian for to be nervous, then the Russian for anxious.
\end{quote}

\subsection*{You'll want to know: \ru{беспокоиться}}
\textbf{\ru{беспокоиться}} (\emph{bespokóitsya}) — \textquotedblleft{}to worry\textquotedblright{}.

\begin{quote}
\ru{Родители} \ru{беспокоятся}, \ru{когда} \ru{я} \ru{поздно} \ru{прихожу} \ru{домой}.
\end{quote}

\begin{quote}
\emph{Rodíteli bespokóyatsya, kogdá ya pózdno prikhozhú domóy.}
\end{quote}

\begin{quote}
My parents worry when I come home late.
\end{quote}

Imperfective. \emph{\ru{Беспокоиться} \ru{о} \ru{ком}} or \emph{\ru{за} \ru{кого}} — to worry about someone. Perfective \ru{побеспокоиться}.

\begin{cousinweb}
\emph{\ru{беспокоиться}} is built on \emph{\ru{без}} 'without' and \emph{\ru{покой}} 'rest, peace', which is inherited from Proto-Slavic \emph{pokoj\ru{ь}} (EtymDB).
\end{cousinweb}

\begin{grammarlens}[title={What you've built}]
One more word for telling, explaining and describing at B1.
\end{grammarlens}

\subsection*{Your turn}
\begin{itemize}
\raggedright
  \item \emph{Say it:} \emph{bespokóitsya}
  \item \emph{Say it:} \emph{bespokóitsya}, once more
  \item \emph{Say it:} a sentence of your own with \emph{bespokóitsya}
  \item \emph{Recall:} say the Russian for to be nervous, then the Russian for anxious, then say \emph{bespokóitsya} again
\end{itemize}

\begin{tcolorbox}[breakable,colback=teal!4,colframe=teal!35!black,title={Before you move on}]
What does \ru{беспокоиться} mean? (\textquotedblleft{}To worry\textquotedblright{}.) Say it once more.
\end{tcolorbox}
\section[\texorpdfstring{\ru{беспокойство} (bespokóystvo)}{беспокойство (bespokóystvo)}]{\ru{беспокойство} (bespokóystvo) — worry, anxiety}

\label{lesson:RU-C474-bespokoystvo}



\begin{quote}
Before the new one: say the Russian for anxious, then the Russian for to worry.
\end{quote}

\subsection*{You'll want to know: \ru{беспокойство}}
\textbf{\ru{беспокойство}} (\emph{bespokóystvo}) — \textquotedblleft{}worry, anxiety\textquotedblright{}.

\begin{quote}
\ru{Его} \ru{беспокойство} \ru{было} \ru{понятно}.
\end{quote}

\begin{quote}
\emph{Yegó bespokóystvo býlo ponyátno.}
\end{quote}

\begin{quote}
His worry was understandable.
\end{quote}

Neuter.

\begin{cousinweb}
\emph{\ru{беспокойство}} is built on \emph{\ru{беспокойный}} 'restless', from \emph{\ru{без}} 'without' and \emph{\ru{покой}} 'rest', inherited from Proto-Slavic \emph{pokoj\ru{ь}} (EtymDB).
\end{cousinweb}

\begin{grammarlens}[title={What you've built}]
One more word for telling, explaining and describing at B1.
\end{grammarlens}

\subsection*{Your turn}
\begin{itemize}
\raggedright
  \item \emph{Say it:} \emph{bespokóystvo}
  \item \emph{Say it:} \emph{bespokóystvo}, once more
  \item \emph{Say it:} a sentence of your own with \emph{bespokóystvo}
  \item \emph{Recall:} say the Russian for anxious, then the Russian for to worry, then say \emph{bespokóystvo} again
\end{itemize}

\begin{tcolorbox}[breakable,colback=teal!4,colframe=teal!35!black,title={Before you move on}]
What does \ru{беспокойство} mean? (\textquotedblleft{}Worry, anxiety\textquotedblright{}.) Say it once more.
\end{tcolorbox}
\section[\texorpdfstring{\ru{успокаиваться} (uspokáivatsya)}{успокаиваться (uspokáivatsya)}]{\ru{успокаиваться} (uspokáivatsya) — to calm down}

\label{lesson:RU-C474-uspokaivatsya}



\begin{quote}
Before the new one: say the Russian for to worry, then the Russian for worry.
\end{quote}

\subsection*{You'll want to know: \ru{успокаиваться}}
\textbf{\ru{успокаиваться}} (\emph{uspokáivatsya}) — \textquotedblleft{}to calm down\textquotedblright{}.

\begin{quote}
\ru{После} \ru{чашки} \ru{чая} \ru{она} \ru{быстро} \ru{успокаивается}.
\end{quote}

\begin{quote}
\emph{Pósle cháshki cháya oná býstro uspokáivayetsya.}
\end{quote}

\begin{quote}
After a cup of tea she soon calms down.
\end{quote}

Imperfective; perfective \ru{успокоиться}. Built like \emph{\ru{спокойный}} — calm.

\begin{cousinweb}
\emph{\ru{успокаиваться}} is the reflexive of \emph{\ru{успокаивать}}, the partner of \emph{\ru{успокоить}} 'to calm', built on \emph{\ru{покой}} 'rest', which is inherited from Proto-Slavic \emph{pokoj\ru{ь}} (EtymDB).
\end{cousinweb}

\begin{grammarlens}[title={What you've built}]
One more word for telling, explaining and describing at B1.
\end{grammarlens}

\subsection*{Your turn}
\begin{itemize}
\raggedright
  \item \emph{Say it:} \emph{uspokáivatsya}
  \item \emph{Say it:} \emph{uspokáivatsya}, once more
  \item \emph{Say it:} a sentence of your own with \emph{uspokáivatsya}
  \item \emph{Recall:} say the Russian for to worry, then the Russian for worry, then say \emph{uspokáivatsya} again
\end{itemize}

\begin{tcolorbox}[breakable,colback=teal!4,colframe=teal!35!black,title={Before you move on}]
What does \ru{успокаиваться} mean? (\textquotedblleft{}To calm down\textquotedblright{}.) Say it once more.
\end{tcolorbox}
